\chapter{El grafo en memoria}
\label{ch:graph-hunter}

El struct central es \Graph{} (con alias \code{GraphHunter} por
compatibilidad). Es el único productor de \code{CompactRelation},
el único dueño del \code{StringInterner} y del
\code{MetadataStore}, y el único target del matcher DFS.

\section{Campos}

\begin{lstlisting}[style=rust]
pub struct InMemoryGraph {
    // Almacenamiento directo.
    pub(crate) entities:         HashMap<StrId, Entity>,
    pub(crate) edge_store:       SpillableEdgeStore,
    pub(crate) interner:         StringInterner,
    pub(crate) meta_store:       MetadataStore,

    // Indices -- desnormalizados para velocidad de consulta.
    pub(crate) reverse_adj:      HashMap<StrId, Vec<StrId>>,
    pub(crate) type_index:       HashMap<EntityType, HashSet<StrId>>,
    pub(crate) entity_type_tags: Vec<u8>, // denso por StrId

    // Internador de datasets (Arc<str>, tag 16-bit).
    pub(crate) dataset_tags:     Vec<Arc<str>>,
    pub(crate) dataset_map:      HashMap<Arc<str>, u16>,
}
\end{lstlisting}

La redundancia entre \code{entities[...].entity\_type} y
\code{entity\_type\_tags[...]} es deliberada: el primero es la
fuente canónica pero cuesta un lookup en HashMap; el segundo es
un array SoA indexado por \code{StrId.0 as usize} --- una lectura
de byte en tiempo constante que el matcher usa dentro de su loop
inner para descartar aristas cuyo tipo de destino no matchea el
paso de la hipótesis.

\begin{invariant}[Coherencia de índices]
Para todo \StrId{} $x$ producido por \code{interner.intern(\_)}:
\code{entities.contains\_key(\&x)} si y sólo si el string fue
agregado como entidad, y en ese caso
\code{type\_index[et].contains(\&x)} con \code{et =
entities[\&x].entity\_type}, y
\code{entity\_type\_tags[x.0 as usize]} es igual al tag numérico
de ese tipo de entidad.
\end{invariant}

\section{Algoritmos de inserción}

\begin{algorithm}
\caption{\textsc{Add-Entity}}\label{alg:add-entity}
\KwIn{grafo $G$, entidad $e = (\text{id}, \text{type})$}
$x \gets \code{G.interner.intern(}\text{e.id}\code{)}$\;
\If{$x \in G.\text{entities}$}{\Return \Comment*[r]{idempotente en dup}}
\code{G.entities.insert(x, e)}\;
\code{G.type\_index[\text{e.type}].insert(x)}\;
\code{G.entity\_type\_tags.resize(max(len, x.0+1), \_)}\;
\code{G.entity\_type\_tags[x.0] $\gets$ tag(\text{e.type})}\;
\code{G.reverse\_adj.entry(x).or\_default()}\;
\Return\;
\end{algorithm}

\begin{algorithm}
\caption{\textsc{Add-Relation}}\label{alg:add-relation}
\KwIn{grafo $G$, relación $r = (\text{src}, \text{dst}, \text{rel}, \tau, M)$}
$s \gets \code{G.interner.get(r.src)}$; fallar si \code{None}\;
$d \gets \code{G.interner.get(r.dst)}$; fallar si \code{None}\;
$o \gets \code{G.meta\_store.append(M)}$\;
$t \gets $ \textsc{Codificar-Ingested-At}(\textsc{Now-Seconds}())\;
$c \gets $ \Compact{}($s, d, \text{rel}, \tau, o, \text{dsTag}, t$)\;
\code{G.edge\_store.append(s, c)}\;
\code{G.reverse\_adj[d].push(s)}\;
\Return\;
\end{algorithm}

\begin{complexity}
\textsc{Add-Entity}: $\Ord{1}$ amortizado en la cantidad de
entidades existentes; el resize de \code{entity\_type\_tags} se
amortiza con la estrategia de duplicación de vectores.
\textsc{Add-Relation}: $\Ord{1}$ amortizado, dominado por el
\code{append} del edge store spillable (Capítulo~\ref{ch:spill})
y el lookup del internador. El espacio por arista insertada es
32 bytes más cualquier byte de metadato apendizado al
\code{MetadataStore}.
\end{complexity}

\section{Finalización}

Las listas de adyacencia dentro de \code{edge\_store} no están
ordenadas por timestamp al insertar --- ordenar en cada
\textsc{Add-Relation} pagaría $\Ord{\log d}$ por inserción y la
ingesta viene de fuentes streaming en bulk. En su lugar, tras un
batch, el llamador invoca \code{sort\_edges\_by\_timestamp}, que
es el disparador externo de \code{SpillableEdgeStore::finalize}:

\begin{algorithm}
\caption{\textsc{Sort-Edges-By-Timestamp}}\label{alg:finalize-sort}
\KwIn{grafo $G$}
\code{G.edge\_store.finalize()}\Comment*[r]{merge k-way, \S\ref{ch:spill}}
\For{$s \in G.\text{entities.keys}$}{
    ordenar \code{G.edge\_store.get\_edges(s)} in place por
    \code{(timestamp, dest\_sid)} \Comment*[r]{\Ord{d(s) \log d(s)}}
}
\end{algorithm}

\begin{complexity}
$\Ord{m \log \bar d}$ donde $\bar d$ es el grado saliente
promedio, ya que cada lista de adyacencia de largo $d(v)$ se
ordena por separado. Si el store estaba en modo spill,
\code{finalize} primero hace un merge $k$-way a costo $\Ord{m \log
k}$ (Capítulo~\ref{ch:spill}). El espacio sólo crece con el
tamaño del archivo de merge durante este paso.
\end{complexity}

\begin{inthecode}
\filepath{graph\_hunter\_core/src/graph.rs:67} --- campos del
struct.\\
\filepath{graph\_hunter\_core/src/graph.rs:211} ---
\textsc{Add-Entity}.\\
\filepath{graph\_hunter\_core/src/graph.rs:227} ---
\textsc{Add-Relation}.\\
\filepath{graph\_hunter\_core/src/graph.rs:437} ---
\textsc{Sort-Edges-By-Timestamp}.
\end{inthecode}
